% Options for packages loaded elsewhere
\PassOptionsToPackage{unicode}{hyperref}
\PassOptionsToPackage{hyphens}{url}
\PassOptionsToPackage{dvipsnames,svgnames,x11names}{xcolor}
%
\documentclass[
  11pt,
]{article}
\usepackage{amsmath,amssymb}
\usepackage{iftex}
\ifPDFTeX
  \usepackage[T1]{fontenc}
  \usepackage[utf8]{inputenc}
  \usepackage{textcomp} % provide euro and other symbols
\else % if luatex or xetex
  \usepackage{unicode-math} % this also loads fontspec
  \defaultfontfeatures{Scale=MatchLowercase}
  \defaultfontfeatures[\rmfamily]{Ligatures=TeX,Scale=1}
\fi
\usepackage{lmodern}
\ifPDFTeX\else
  % xetex/luatex font selection
\fi
% Use upquote if available, for straight quotes in verbatim environments
\IfFileExists{upquote.sty}{\usepackage{upquote}}{}
\IfFileExists{microtype.sty}{% use microtype if available
  \usepackage[]{microtype}
  \UseMicrotypeSet[protrusion]{basicmath} % disable protrusion for tt fonts
}{}
\makeatletter
\@ifundefined{KOMAClassName}{% if non-KOMA class
  \IfFileExists{parskip.sty}{%
    \usepackage{parskip}
  }{% else
    \setlength{\parindent}{0pt}
    \setlength{\parskip}{6pt plus 2pt minus 1pt}}
}{% if KOMA class
  \KOMAoptions{parskip=half}}
\makeatother
\usepackage{xcolor}
\usepackage[margin=1in]{geometry}
\usepackage{longtable,booktabs,array}
\usepackage{calc} % for calculating minipage widths
% Correct order of tables after \paragraph or \subparagraph
\usepackage{etoolbox}
\makeatletter
\patchcmd\longtable{\par}{\if@noskipsec\mbox{}\fi\par}{}{}
\makeatother
% Allow footnotes in longtable head/foot
\IfFileExists{footnotehyper.sty}{\usepackage{footnotehyper}}{\usepackage{footnote}}
\makesavenoteenv{longtable}
\setlength{\emergencystretch}{3em} % prevent overfull lines
\providecommand{\tightlist}{%
  \setlength{\itemsep}{0pt}\setlength{\parskip}{0pt}}
\setcounter{secnumdepth}{-\maxdimen} % remove section numbering
\ifLuaTeX
  \usepackage{selnolig}  % disable illegal ligatures
\fi
\IfFileExists{bookmark.sty}{\usepackage{bookmark}}{\usepackage{hyperref}}
\IfFileExists{xurl.sty}{\usepackage{xurl}}{} % add URL line breaks if available
\urlstyle{same}
\hypersetup{
  pdfauthor={Stefano Alimonti --- stefalimonti@gmail.com},
  colorlinks=true,
  linkcolor={blue},
  filecolor={Maroon},
  citecolor={Blue},
  urlcolor={blue},
  pdfcreator={LaTeX via pandoc}}

\author{Stefano Alimonti --- stefalimonti@gmail.com}
\date{March 2026}

\begin{document}

{
\hypersetup{linkcolor=black}
\setcounter{tocdepth}{2}
\tableofcontents
}
\section{The Sector Ratio in Binary Multiplication: From Markov Failure to
Transcendence}\label{the-sector-ratio-in-binary-multiplication-from-markov-failure-to-transcendence}

\textbf{Author:} Stefano Alimonti \textbf{Affiliation:} Independent Researcher
\textbf{Date:} March 2026

\begin{center}\rule{0.5\linewidth}{0.5pt}\end{center}

\newpage

\subsection{Abstract}\label{abstract}

We study the sector partition function of binary multiplication: given K-bit
integers X, Y with product P = XY of D = 2K−1 bits, we partition the
4\^{}\{K−1\} digit pairs by the second-highest bits (x\_\{K−2\}, y\_\{K−2\}) and
compute the carry-weighted ratio R(K) = σ₁₀/σ₀₀ under two natural carry
valuations. The \emph{schoolbook} valuation uses the carry at the fixed position
D − 2; the \emph{cascade} valuation uses a first-passage stopping time (the
carry one step below the topmost nonzero carry). The Markov model of independent
carry propagation predicts R\_Markov = +2/3. We show three independent failures
of this prediction:

\begin{enumerate}
\def\labelenumi{\arabic{enumi}.}
\tightlist
\item
  The Markov model gives the \emph{wrong sign} (+2/3 vs.~the negative R(K)
  observed for K ≥ 7).
\item
  It misses a qualitative ``J = 1'' exclusion effect that removes the dominant
  contribution.
\item
  The digit correlations follow the exact profile \(\rho(d) = (K-d)/(3K+2)\),
  with leading-order Fejér approximation \(\rho(d) \approx (K-d)/(3K)\), and
  couple the Dirichlet modes of the carry bridge.
\end{enumerate}

Exact enumeration up to K = 21 shows R\^{}\{cas\}(K) → −π with asymptotic
convergence rate \(O(K^2 \cdot 2^{-K})\) {[}P1, §9.3{]}, while the schoolbook
limit R\^{}\{sb\}(∞) = R₀ = −3.9312\ldots{} is a purely logarithmic constant (no
π). The decomposition R\^{}\{cas\} = R₀ + ΔR isolates all transcendental content
in the cascade correction ΔR = +0.7896\ldots{} {[}P1, §10{]}. We present the
full computational pipeline, analyze the digit correlation structure, and
quantify the spectral mechanism by which rational arithmetic produces a
transcendental constant.

When the Fejér coupling is treated as a perturbation of the Markov operator, the
effective coupling diverges to the pole A = 3/2 of the Möbius transformation ---
a universal phenomenon across all bases. The physical system avoids this
divergence through non-linear carry dynamics, operating at A* = 3/2 −
1/(2(π+1)).

We prove that no D-odd pair has sector (1,1), that the count ratio N₁₀/N₀₀
converges to the closed form (2ln(4/3) − 1/2)/(2ln(9/8)) = 0.31993\ldots{}
(involving only ln 2 and ln 3), and that the entire factor π in the conjectured
limit R → −π enters exclusively through the carry weight ratio ⟨val⟩₁₀/⟨val⟩₀₀,
not the pair counts. An extended Markov chain with W bits of memory demonstrates
that bit-sharing correlations between adjacent convolutions are the mechanism
transforming the Markov super-exponential decay Q ∼ (3/4)\^{}d into the
empirical geometric decay Q ≈ 0.55, with the effective per-step survival
probability converging toward 1/2 --- the Diaconis--Fulman eigenvalue.

\begin{center}\rule{0.5\linewidth}{0.5pt}\end{center}

\newpage

\subsection{1. Introduction}\label{introduction}

This is the computational companion to {[}P1{]}. Where that paper develops the
analytical theory (the closed-form Möbius transformation, the spectral zeta
function, the Dirichlet character connection, the R₀ + ΔR decomposition, and the
Bernoulli hierarchy), this paper provides:

\begin{enumerate}
\def\labelenumi{\arabic{enumi}.}
\tightlist
\item
  The complete definition and algorithmic computation of R(K) under both the
  schoolbook and cascade valuations (Section 2).
\item
  A demonstration that the Markov model fails \emph{qualitatively} (wrong sign,
  wrong magnitude, wrong mechanism).
\item
  The extraction and characterization of digit correlations (the Fejér kernel).
\item
  The spectral consequences of the coupling, including the universal pole
  divergence.
\item
  Exact enumeration data through K = 21.
\item
  Structural decomposition: the top-carry constraint, closed-form count ratio,
  and the isolation of π to carry weights (Section 9).
\item
  The bit-sharing mechanism: an extended Markov chain with W bits of memory that
  reproduces Q\_emp ≈ 0.55 and connects the cascade stopping time to the
  Diaconis--Fulman eigenvalue (Section 9.5).
\end{enumerate}

For the angular uniqueness of base 2 (why π enters only in binary), see {[}G{]}.

\subsubsection{1.1 Motivation}\label{motivation}

The question is elementary: in the schoolbook multiplication of two random K-bit
numbers, how does the carry at position D − 2 depend on which \emph{sector} the
input falls in? The answer turns out to involve π, emerging through a mechanism
that has no classical geometric content.

\begin{center}\rule{0.5\linewidth}{0.5pt}\end{center}

\newpage

\subsection{2. Definitions and Setup}\label{definitions-and-setup}

\subsubsection{2.1 Binary multiplication}\label{binary-multiplication}

Let X = Σ\_\{i=0\}\^{}\{K-1\} x\_i 2\^{}i and Y = Σ\_\{i=0\}\^{}\{K-1\} y\_i
2\^{}i with x\_\{K−1\} = y\_\{K−1\} = 1 (K-bit integers). The product P = XY has
D = 2K − 1 bits. At each position j:

\begin{verbatim}
conv_j = Σ_{i+i'=j, 0≤i,i'≤K-1} x_i · y_{i'}

p_j + 2·c_{j+1} = conv_j + c_j,    c_0 = 0
\end{verbatim}

\textbf{Definition 1a (Schoolbook carry weight).} For a pair (X, Y), define the
\emph{schoolbook weight} w\^{}\{sb\}(X,Y) = 2c\_\{D−2\} − 1 ∈ \{−1, +1\}, using
the carry at the fixed position D − 2.

\textbf{Definition 1b (Cascade carry weight).} Let M(X,Y) = max\{j : c\_j ≥ 1\}
be the position of the topmost nonzero carry (the \emph{cascade stopping
depth}). The \emph{cascade weight} is w\^{}\{cas\}(X,Y) = c\_\{M−1\} − 1 ∈ \{−1,
0, +1\}, the carry one step below the stopping point. For pairs with M = 0 we
set w\^{}\{cas\} = 0.

\textbf{Remark (Two valuations).} The schoolbook weight uses a fixed observation
point (position D − 2), while the cascade weight uses a random stopping time
(position M − 1). The two coincide when M = D − 1 (the carry chain extends to
the top) and differ otherwise. Approximately 55\% of sector-(0,0) pairs have M =
D − 1; for sector (1,0), the top-carry constraint (Proposition 6) forces M ≤ D −
2 for all pairs. The distinction is analysed fully in {[}P1, §10{]}.

\textbf{Definition 2 (Sector partition function).} For a, c ∈ \{0, 1\} and
valuation ∗ ∈ \{sb, cas\}:

\begin{verbatim}
σ^∗_ac(K) = Σ_{(X,Y): x_{K-2}=a, y_{K-2}=c} w^∗(X,Y)
\end{verbatim}

\textbf{Definition 3 (Sector ratio).} R\^{}∗(K) =
σ\textsuperscript{∗₁₀(K)/σ}∗₀₀(K).

The cascade ratio R\^{}\{cas\}(K) → −π (Conjecture 1). The schoolbook ratio
R\^{}\{sb\}(K) → R₀ = (1/2 − 2ln(4/3))/(1 + ln(3/8)) = −3.9312\ldots, a purely
logarithmic constant {[}E, Proposition 1{]}. The decomposition R\^{}\{cas\} = R₀
+ ΔR with ΔR = +0.7896\ldots{} isolates all transcendental content in the
cascade correction; see {[}P1, §10.1{]}.

\begin{center}\rule{0.5\linewidth}{0.5pt}\end{center}

\newpage

\subsection{3. The Markov Model and Its
Failure}\label{the-markov-model-and-its-failure}

\subsubsection{3.1 The independent-convolution
approximation}\label{the-independent-convolution-approximation}

The Diaconis--Fulman Markov chain {[}1, 2, 4{]} models the carry sequence (c\_0,
c\_1, \ldots, c\_\{D−1\}) as a Markov chain: the carry into position j+1 depends
only on c\_j and the convolution conv\_j, which is treated as an independent
random variable given only j (not the actual digits).

\subsubsection{3.2 Qualitative failure: wrong
sign}\label{qualitative-failure-wrong-sign}

\textbf{Proposition 1} (computational). The Markov model predicts R\_Markov(K)
\textgreater{} 0 for all K, while exact enumeration gives R(K) \textless{} 0 for
K ≥ 7.

\emph{Evidence.} Verified by exact enumeration for K = 5, \ldots, 21 and
independently confirmed at K = 20, 21 by the C-language enumerator (E44). A
formal proof would require showing that the Markov per-sector weights have
opposite sign to the exact per-sector weights; the mechanism (J = 1 exclusion,
Proposition 2) provides the structural explanation.

This is not a quantitative error --- the Markov model predicts the \emph{wrong
direction} of the sector effect. The sign flip is the most dramatic signature of
non-Markovian behavior in the carry chain.

\begin{longtable}[]{@{}lll@{}}
\toprule\noalign{}
K & R\_Markov(K) & R\^{}\{cas\}(K) \\
\midrule\noalign{}
\endhead
\bottomrule\noalign{}
\endlastfoot
5 & +0.59 & +0.25 \\
7 & +0.62 & −0.09 \\
9 & +0.64 & −0.73 \\
11 & +0.65 & −1.554 \\
13 & +0.65 & −2.339 \\
15 & +0.66 & −2.824 \\
17 & +0.66 & −3.035 \\
19 & +0.66 & −3.110 \\
20 & +0.66 & −3.125 \\
21 & +0.66 & −3.133 \\
\end{longtable}

\subsubsection{3.3 The J = 1 exclusion effect}\label{the-j-1-exclusion-effect}

The Markov model's dominant contribution to σ₁₀ comes from the ``J = 1 pathway''
--- a specific carry pattern near the top of the bridge. In the exact
computation, this pathway is \emph{forbidden} by a digit constraint that the
Markov model ignores.

\textbf{Proposition 2 (J = 1 exclusion, computationally verified).} In the
sector (a, c) = (1, 0), the leading-order Markov contribution at position J = 1
(one step from the boundary) requires digit configurations that violate the
D-odd constraint. Removing this contribution flips the sign of σ₁₀.

\emph{Evidence.} Verified by direct comparison of the Markov and exact
computations for K = 5, \ldots, 13 (experiment E06). The formal case analysis on
digit configurations at positions K−2 and K−1 is outlined but not fully proved;
see Open Problem 3.

\begin{center}\rule{0.5\linewidth}{0.5pt}\end{center}

\newpage

\subsection{4. Digit Correlations}\label{digit-correlations}

\subsubsection{4.1 Inter-position
correlations}\label{inter-position-correlations}

In the Markov model, convolutions at positions j and j' are independent given
the carry state. In reality, the shared digit structure induces correlations.

\textbf{Definition 4 (Digit-position correlation).} For positions j, j' with
\textbar j − j'\textbar{} = d:

\begin{verbatim}
ρ(d) = Cov[conv_j, conv_{j'}] / Var[conv_j]
\end{verbatim}

\subsubsection{4.2 Correlation profile: the Fejér
kernel}\label{correlation-profile-the-fejuxe9r-kernel}

\textbf{Theorem 6 (Correlation profile).} For the central region of the product
(positions j ≈ K), with K-bit uniformly random inputs:

\begin{verbatim}
ρ(d) = (K − d)/(3K + 2),    0 ≤ d ≤ K−1
\end{verbatim}

(exact, from Var{[}conv\_j{]} = (3K+2)/16 and Cov = (K−d)/16). The Fejér kernel
approximation \(\rho(d) \approx (K-d)/(3K)\) holds with absolute error
\(2(K-d)/(3K(3K+2)) = O(1/K)\) uniformly in \(d\); it becomes exact in the limit
\(K \to \infty\) at fixed \(d/K\).

\emph{Proof.} The convolution conv\_j = Σ\_\{i+i'=j\} x\_i y\_\{i'\} is a sum of
products of independent bits. Two convolutions conv\_j and conv\_\{j+d\} share
exactly K − d digit pairs (when both indices are in the central region). Each
shared pair (x\_i, y\_\{j−i\}) = (x\_i, y\_\{j+d−i\}) contributes Cov(x\_i
y\_\{j−i\}, x\_i y\_\{j+d−i\}) = E{[}x\_i² y\_\{j−i\} y\_\{j+d−i\}{]} − E{[}x\_i
y\_\{j−i\}{]} E{[}x\_i y\_\{j+d−i\}{]}. Since x\_i ∈ \{0,1\}, x\_i² = x\_i, so
E{[}x\_i² y\_\{j−i\} y\_\{j+d−i\}{]} = (1/2)(1/2)(1/2) = 1/8 (for interior terms
where all factors are random), while E{[}x\_i y\_\{j−i\}{]} E{[}x\_i
y\_\{j+d−i\}{]} = (1/4)² = 1/16. Each shared pair contributes 1/8 − 1/16 = 1/16.
The total covariance is (K − d)/16. The variance of conv\_j in the central
region has two regimes: the K − 2 interior terms have Var{[}x\_i y\_\{j−i\}{]} =
E{[}(x\_i y\_\{j−i\})\^{}2{]} − (E{[}x\_i y\_\{j−i\}{]})\^{}2 = 1/4 − 1/16 =
3/16 each, and the 2 border terms (involving the fixed MSBs) have Var = 1/4
each. Since terms in the same convolution are independent (each uses a distinct
pair of bit indices), Var{[}conv\_j{]} = (K − 2) · 3/16 + 2 · 1/4 = (3K + 2)/16
≈ 3K/16 for K ≫ 1. Hence ρ(d) = Cov/Var = (K−d)/16 ÷ (3K+2)/16 = (K−d)/(3K+2). □

\subsubsection{4.3 Implications}\label{implications}

The Fejér kernel has three notable properties: - Triangular (linear decay with
distance) - Positive everywhere (no anti-correlations) - Width proportional to K
(correlations span the entire bridge)

This long-range correlation is what pushes the effective eigenvalue shift past
the critical point A = 1, converting the rational Markov result into the
transcendental −π {[}P1, Section 5{]}.

\begin{center}\rule{0.5\linewidth}{0.5pt}\end{center}

\newpage

\subsection{5. The Exact Transfer Operator}\label{the-exact-transfer-operator}

\subsubsection{5.1 State space}\label{state-space}

The exact (non-Markov) transfer operator tracks both the carry \emph{and} the
last d digits of both X and Y (the ``digit buffer''). For buffer depth d, the
state space has \textbar S\textbar{} = 2 × 2\^{}d × 2\^{}d = 2\^{}\{2d+1\}
states.

\subsubsection{5.2 Eigenvalue spectrum}\label{eigenvalue-spectrum}

\begin{longtable}[]{@{}llll@{}}
\toprule\noalign{}
Buffer depth d & \textbar S\textbar{} & Largest eigenvalue & Next eigenvalue \\
\midrule\noalign{}
\endhead
\bottomrule\noalign{}
\endlastfoot
0 (Markov) & 2 & 1/2 & −1/2 \\
1 & 8 & 1/2 & 1/2 \\
2 & 32 & 0.3536 & 0.3536 \\
3 & 128 & 0.3536 & 0.3536 \\
4 & 512 & 0.3460 & 0.3460 \\
\end{longtable}

The eigenvalue \(0.3536 \approx 1/\sqrt{8} = 1/(2\sqrt{2})\) at levels 2--3 is
consistent with a \(1/b^{3/2}\) scaling for the first non-Markov mode in base 2.
The exact spectrum is much richer than the Markov model's two eigenvalues
\{+1/2, −1/2\}. Most additional eigenvalues are parasitic (related to
digit-buffer dynamics, not carry propagation), but they collectively produce the
correlation-induced shift that generates −π.

\subsubsection{5.3 Non-universality of the exact
spectrum}\label{non-universality-of-the-exact-spectrum}

While the Markov spectral zeta function converges to ζ(2s) {[}P1, Theorem 3{]},
the exact spectral zeta function does not --- the parasitic eigenvalues from the
digit buffer break the Weyl asymptotics. The conjectured −π limit comes from a
\emph{projected} quantity (the sector ratio) that isolates the carry degrees of
freedom.

\begin{center}\rule{0.5\linewidth}{0.5pt}\end{center}

\newpage

\subsection{6. The Fejér Kernel and Dirichlet
Coupling}\label{the-fejuxe9r-kernel-and-dirichlet-coupling}

\subsubsection{6.1 Coupling matrix}\label{coupling-matrix}

Using the leading-order Fejér approximation \(\rho(d) \approx (K-d)/(3K)\)
(Theorem 6; exact form \(\rho(d) = (K-d)/(3K+2)\)), the digit-correlation
profile couples the Dirichlet modes φ\_n(j) = sin(nπj/L) through

\begin{verbatim}
V_{nn'} = Σ_{d=0}^{K-1} ρ(d) Σ_j φ_n(j) φ_{n'}(j+d)
\end{verbatim}

\textbf{Proposition 3 (Non-uniform diagonal coupling; numerical).} The diagonal
elements V\_\{nn\} are not proportional to sin²(nπ/(2L)) for all n.~The
effective shift A\_eff(n) = V\_\{nn\}/sin²(nπ/(2L)) varies from
\textasciitilde400 (low modes) to \textasciitilde0.2 (high modes) at K = 15.

Verified numerically in experiment E10; no closed-form proof is given.

\subsubsection{6.2 Universality argument}\label{universality-argument}

Despite the non-uniformity of A\_eff(n), the sector ratio in the bridge model is
determined by a single number --- the effective A parameter in the Möbius
formula S(A) = 2(1−A)/(3−2A) {[}P1, Corollary 1{]}. The mapping from the full
coupling matrix V\_\{nn'\} to the effective A involves a resummation that
averages over mode numbers.

\subsubsection{6.3 The bare coupling hits the
pole}\label{the-bare-coupling-hits-the-pole}

The full spectral sum S\_full = v\^{}T (I − D − V)\^{}\{-1\} e, where D is the
Markov diagonal and V the Fejér coupling matrix, was computed numerically:

\begin{longtable}[]{@{}lll@{}}
\toprule\noalign{}
K & S\_full & A\_eff (from Möbius inversion) \\
\midrule\noalign{}
\endhead
\bottomrule\noalign{}
\endlastfoot
10 & +56.0 & 1.509 \\
20 & −8.9 × 10⁶ & 1.500000 \\
50 & +3.3 × 10⁹ & 1.500000 \\
100 & +2.3 × 10⁹ & 1.500000 \\
\end{longtable}

\textbf{Finding 1.} A\_eff → 3/2 exactly --- the pole of the Möbius
transformation. This is \emph{structurally inevitable}: the Fejér kernel has
total strength ρ̂(0) = (K+1)/6 → ∞, and the diagonal V\_nn/sin²(nπ/(2L)) grows as
\textasciitilde K³ for low modes. Any coupling that makes S diverge maps to A =
3/2 via Möbius inversion, regardless of the coupling's structure.

\textbf{Finding 2 (Base universality).} The divergence A\_eff → 3/2 occurs for
all bases tested (b = 2, 3, 5). This confirms that the pole-hitting is a
structural property of the Fejér coupling, not a base-2 artifact.

\subsubsection{6.4 The χ₄ character from midpoint
geometry}\label{the-ux3c7ux2084-character-from-midpoint-geometry}

The Dirichlet character χ₄(n) = sin(nπ/2) arises geometrically. The sector
perturbation is at position j* = K − 2 on a bridge of length L = 2K − 1. Since

\begin{verbatim}
j*/L = (K−2)/(2K−1) → 1/2    as K → ∞,
\end{verbatim}

the Dirichlet mode evaluation gives sin(nπj\emph{/L) → sin(nπ/2) = χ₄(n). The
sector perturbation naturally selects the midpoint of the bridge, and the
midpoint evaluation of the Dirichlet modes }is* the character χ₄. This geometric
origin explains why the sector ratio involves Leibniz's series and L(1, χ₄) =
π/4.

\begin{center}\rule{0.5\linewidth}{0.5pt}\end{center}

\newpage

\subsection{7. Computational Pipeline}\label{computational-pipeline}

\subsubsection{7.1 Algorithm}\label{algorithm}

For each K from 3 to 21:

\begin{enumerate}
\def\labelenumi{\arabic{enumi}.}
\tightlist
\item
  Enumerate all (X, Y) ∈ {[}2\^{}\{K−1\}, 2\^{}K)² (total: 4\^{}\{K−1\} pairs).
\item
  For each pair, compute the full carry chain c\_0, c\_1, \ldots, c\_\{D−1\} by
  schoolbook multiplication {[}5{]}.
\item
  Classify by sector (a, c) = (x\_\{K−2\}, y\_\{K−2\}).
\item
  Compute both valuations: the schoolbook weight w\^{}\{sb\} = 2c\_\{D−2\} − 1
  and the cascade weight w\^{}\{cas\} = c\_\{M−1\} − 1, where M = max\{j : c\_j
  ≥ 1\}.
\item
  Accumulate into σ\^{}\{sb\}\_ac and σ\^{}\{cas\}\_ac.
\item
  Compute R\^{}\{sb\}(K) and R\^{}\{cas\}(K).
\end{enumerate}

Unless otherwise stated, all tables in this paper report the cascade valuation
R\^{}\{cas\}(K), which converges numerically toward −π (Conjecture 1 of
{[}P1{]}).

\subsubsection{7.2 Complexity and
implementation}\label{complexity-and-implementation}

The enumeration requires O(K · 4\^{}K) operations. For K = 20, this is
\textasciitilde5.5 × 10¹² operations over 2.75 × 10¹¹ pairs, of which 1.06 ×
10¹¹ satisfy the D-odd constraint.

Exact enumeration covers K ≤ 21. The K = 20 computation uses a C implementation
with OpenMP 8-thread parallelism, completing in 3344 seconds (\textasciitilde56
minutes) on a modern workstation. Memory usage is O(1) (streaming enumeration,
no storage of pairs).

\textbf{K = 20 exact data (cascade valuation):} σ\^{}\{cas\}₀₀ = −1,967,962,747;
σ\^{}\{cas\}₁₀ = 6,149,608,524; σ₀₀\^{}\{c=0\} = −1,683,746,513; σ₀₀\^{}\{c=1\}
= −284,216,234. R\^{}\{cas\}(20) = −3.12486023; α =
σ₀₀\textsuperscript{\{c=1\}/σ₀₀}\{c=0\} = 0.16880. At K = 21: R\^{}\{cas\}(21) =
−3.13340.

\subsubsection{7.3 Convergence to −π}\label{convergence-to-ux3c0}

All convergence data below use the cascade valuation R\^{}\{cas\}(K). For the
schoolbook valuation, see {[}E, Proposition 1{]} where R\^{}\{sb\}(∞) = R₀ =
−3.9312\ldots{} is proved analytically.

\textbf{Proposition 4 (Richardson working ansatz; numerical).} The empirical gap
ratios \(g(K)/g(K-1)\) converge monotonically to \(1/2\), consistent with an
exponential ansatz

\[R^{cas}(K) = R^{cas}(\infty) + c_1 \rho^K + c_2 \rho^{2K} + \cdots, \qquad \rho = 1/2.\]

The Neville--Aitken table with step-size \(h = \rho^K\) eliminates successive
error terms.

\emph{Note.} A polynomial ansatz in \(1/K^2\) was initially considered but
diverges under Neville extrapolation. The exponential form reflects the true
asymptotic rate \(O(K^2 \cdot 2^{-K})\) {[}P1, §9.3{]}.

\begin{longtable}[]{@{}llll@{}}
\toprule\noalign{}
K & R(K) & ε(K) = \textbar R(K) + π\textbar{} & Sig. digits \\
\midrule\noalign{}
\endhead
\bottomrule\noalign{}
\endlastfoot
11 & −1.554 & 1.588 & --- \\
13 & −2.339 & 0.803 & 0.1 \\
15 & −2.824 & 0.317 & 0.5 \\
17 & −3.035 & 0.107 & 1.0 \\
19 & −3.110 & 0.032 & 1.5 \\
20 & −3.125 & 0.017 & 1.8 \\
21 & −3.133 & 0.008 & 2.1 \\
∞ (Richardson) & −3.1419\ldots{} & --- & 4.0 \\
\end{longtable}

\textbf{Proposition 5 (Richardson extrapolation; numerical).} Applying
Neville--Aitken acceleration {[}6{]} with the exponential ansatz of Proposition
4 to the subsequence R(7), R(9), \ldots, R(21) yields R\_extrap =
−3.1419\ldots{} to ≈ 4.0 significant digits (4.4 when restricted to exact
values). The extrapolation is stable under cross-validation: removing any single
K ∈ \{7, \ldots, 17\} changes the estimate by less than 10⁻⁴. This is a
numerical observation, not a convergence proof.

\begin{center}\rule{0.5\linewidth}{0.5pt}\end{center}

\newpage

\subsection{8. Physical Analogies}\label{physical-analogies}

The sign flip from +2/3 (Markov) to −π (exact) has structural parallels in
mathematical physics. In quantum field theory, a classically scale-invariant
theory acquires a non-zero trace of the energy-momentum tensor through quantum
corrections --- here the ``classical'' (Markov) result is rational and the
``quantum'' (correlation-corrected) result is transcendental, with summing over
Dirichlet eigenmodes providing the mechanism. The carry bridge also shares the
spectral structure of a 1D Ising model with fixed-spin boundaries: Dirichlet
eigenfunctions, exponential decay with correlation length ξ = 1/log(1/λ₂), and
π-dependence through the spectral decomposition.

These analogies are structural, not rigorous equivalences. The ``bare'' Fejér
coupling diverges (Section 6.3), while the physical system operates at a finite
effective coupling A* = 3/2 − 1/(2(π+1)). Formalizing this renormalization
mechanism remains open.

\begin{center}\rule{0.5\linewidth}{0.5pt}\end{center}

\newpage

\subsection{9. Structural Decomposition of R
(E16--E19)}\label{structural-decomposition-of-r-e16e19}

\subsubsection{9.1 The top-carry constraint}\label{the-top-carry-constraint}

\textbf{Proposition 6 (Top-carry identity).} conv\_\{D−2\} = a + c for all D-odd
pairs, where a = x\_\{K−2\}, c = y\_\{K−2\} are the sector bits.

\emph{Proof.} conv\_\{D−2\} = Σ\_\{i+i'=D−2\} x\_i y\_\{i'\} with i, i' ∈ {[}0,
K−1{]}. The only solutions are (i, i') = (K−2, K−1) and (K−1, K−2), giving
conv\_\{D−2\} = x\_\{K−2\}·y\_\{K−1\} + x\_\{K−1\}·y\_\{K−2\} = a·1 + 1·c = a +
c.~□

\textbf{Corollary (Sector constraints).} - Sector (1,1): conv\_\{D−2\} = 2,
D-odd requires carries{[}D−1{]} = 0, but ⌊(2 + carries{[}D−2{]})/2⌋ ≥ 1.
Contradiction → \textbf{N₁₁ = 0 for all K}. - Sector (1,0): conv\_\{D−2\} = 1,
so carries{[}D−2{]} must be 0 → \textbf{M ≤ D−3} forced. - Sector (0,0):
conv\_\{D−2\} = 0, carries{[}D−2{]} ∈ \{0,1\} → two sub-populations: ``top'' (M
= D−2, \textasciitilde55\%) and ``low'' (M \textless{} D−2).

This gives σ₀₀ = σ₀₀\^{}\{top\} + σ₀₀\^{}\{low\} and σ₁₀ = σ₁₀\^{}\{low\} only.

\subsubsection{9.2 Closed-form count ratio}\label{closed-form-count-ratio}

\textbf{Theorem 7 (Count ratio).} In the continuous limit K → ∞:

\begin{verbatim}
N₁₀/N₀₀ → (2ln(4/3) − 1/2) / (2ln(9/8)) = 0.31992784225...
\end{verbatim}

with correction O(1/2\^{}K). The sector fractions of D-odd pairs are:

\begin{longtable}[]{@{}lll@{}}
\toprule\noalign{}
Sector & Fraction of D-odd pairs & Value \\
\midrule\noalign{}
\endhead
\bottomrule\noalign{}
\endlastfoot
(0,0) & 2 ln(9/8) / (2 ln 2 − 1) & 0.6098 \\
(1,0) & (2 ln(4/3) − 1/2) / (2 ln 2 − 1) & 0.1951 \\
(0,1) & same as (1,0) & 0.1951 \\
(1,1) & 0 & 0 \\
\end{longtable}

\emph{Proof.} Write X = 2\^{}\{K−1\}(1+u), Y = 2\^{}\{K−1\}(1+v). D-odd:
(1+u)(1+v) \textless{} 2. Sectors from ⌊2u⌋, ⌊2v⌋. Elementary integration. □

Verified to 8 significant digits against exact enumeration (K = 3--19) with
geometric Richardson extrapolation.

\subsubsection{9.3 The π isolation}\label{the-ux3c0-isolation}

Since R = (⟨val⟩₁₀/⟨val⟩₀₀) × (N₁₀/N₀₀) and N₁₀/N₀₀ involves only ln 2 and ln 3:

\begin{verbatim}
⟨val⟩₁₀/⟨val⟩₀₀ → −π · 2 ln(9/8) / (2 ln(4/3) − 1/2) = −9.8197...
\end{verbatim}

\textbf{The entire factor π in the conjectured limit R → −π enters through the
carry weight ratio}, not through the pair counts. This reduces the proof of
Conjecture 1 to proving this single carry weight identity.

\subsubsection{9.4 The boundary layer mechanism
(E18)}\label{the-boundary-layer-mechanism-e18}

The total variation distance TV(j) = (1/2) Σ\_v \textbar P(c\_j\textbar10) −
P(c\_j\textbar00)\textbar{} reveals a boundary layer with two peaks: - j = K−2
(sector bit): TV ≈ 0.16, from the direct convolution perturbation. - j = D−2
(top constraint): TV ≈ 0.54, from the structural constraint conv\_\{D−2\} = a+c.

The c\_\{M−1\} distribution shifts dramatically:

\begin{longtable}[]{@{}llll@{}}
\toprule\noalign{}
& c\_\{M−1\} = 0 (val = −1) & c\_\{M−1\} = 1 (val = 0) & c\_\{M−1\} = 2 (val =
+1) \\
\midrule\noalign{}
\endhead
\bottomrule\noalign{}
\endlastfoot
Sector (0,0)\_low & 11.2\% & 82.2\% & 6.6\% \\
Sector (1,0) & 2.1\% & 67.9\% & 30.0\% \\
\end{longtable}

This 6--7× shift in the val = +1 proportion drives σ₁₀ \textgreater{} 0 while
σ₀₀ \textless{} 0.

\subsubsection{9.5 The bit-sharing mechanism
(E31)}\label{the-bit-sharing-mechanism-e31}

The Markov model treats convolutions at successive depths as independent given
the carry state, predicting a survival probability Q\_Markov(d) that decays
super-exponentially: Q ∼ (3/4)\^{}d.~Exact enumeration shows geometric decay
Q\_emp ≈ 0.55, an order-of-magnitude discrepancy at moderate depths. The
mechanism is \emph{bit-sharing}: the digit x\_\{K−3\} (for instance) contributes
to convolutions at depths d = 1 and d = 2, inducing correlations that the Markov
model ignores.

\textbf{Definition 5 (Extended Markov chain).} For window size W ≥ 0, define a
Markov chain whose state includes the carry value and a buffer of W bits of both
X and Y. The convolution at each depth d ≤ W is computed exactly from the known
bits; for d \textgreater{} W, the unknown bits are modeled as independent
Bernoulli(1/2). The entry carry at depth W is forward-propagated from c = 0
through the exact chain.

\textbf{Finding 3 (Bit-sharing IS the mechanism).} The extended model with W =
8--10 reproduces the empirical survival probabilities to within 2\%:

\begin{longtable}[]{@{}lllll@{}}
\toprule\noalign{}
d & Q\_Markov & W = 8 & W = 10 & Empirical (K = 10) \\
\midrule\noalign{}
\endhead
\bottomrule\noalign{}
\endlastfoot
2 & 0.500 & 0.542 & 0.543 & 0.592 \\
3 & 0.406 & 0.576 & 0.578 & 0.568 \\
4 & 0.328 & 0.534 & 0.540 & 0.588 \\
5 & 0.264 & 0.522 & 0.533 & 0.627 \\
6 & 0.211 & 0.491 & 0.518 & 0.699 \\
\end{longtable}

At d = 3, the W = 10 model matches the K = 10 empirical value to 1.8\%. The Q
values converge monotonically with W, confirming that each additional bit of
memory captures more of the inter-depth correlation.

\textbf{Finding 4 (Effective eigenvalue → 1/2).} The geometric-fit effective
eigenvalue q₀(W) (defined by Q(d) ≈ q₀ for large d within the window) decreases
toward 1/2 as W grows:

\begin{longtable}[]{@{}ll@{}}
\toprule\noalign{}
W & q₀ \\
\midrule\noalign{}
\endhead
\bottomrule\noalign{}
\endlastfoot
3 & 0.543 \\
5 & 0.515 \\
7 & 0.499 \\
10 & 0.483 \\
\end{longtable}

The limit q₀ → 1/2 = \textbar B₁\textbar{} would recover the Diaconis--Fulman
spectral gap {[}1{]} as the universal per-step survival probability of the
cascade, consistent with the Bernoulli hierarchy {[}P1, §8{]}.

\textbf{Remark.} The extended model successfully explains \emph{why} Q ≈ 0.55
but does not directly produce R\^{}\{cas\} → −π: the denominator σ\^{}\{cas\}₀₀
passes through zero as W increases, causing numerical instability. The
convergence R → −π requires the full carry chain (all K bits), not just the
boundary layer window. The analytical proof of ΔR = −π − R₀ remains open {[}P1,
§11, item 1{]}.

\begin{center}\rule{0.5\linewidth}{0.5pt}\end{center}

\newpage

\subsection{10. Open Problems}\label{open-problems}

\begin{enumerate}
\def\labelenumi{\arabic{enumi}.}
\item
  \textbf{Prove R\^{}\{cas\} = −π analytically.} The decomposition R\^{}\{cas\}
  = R₀ + ΔR (with R₀ proved in {[}E, Proposition 1{]}) reduces the problem to
  showing ΔR = −π − R₀ = +0.7896\ldots. Five equivalent formulations are
  developed in {[}P1, §10.5{]}: series identity over depths, angular integral,
  carry weight computation, resolvent trace, and separable weight matrix with
  carry constraints. The most promising route is the stopping-time series {[}P1,
  §8.2a{]}, which decomposes R through per-depth universal constants Δ(τ). The
  exact analytical mechanism underlying the convergence is identified
  conditionally in {[}E{]} via \textbf{resolvent universality} of the effective
  non-Markovian transfer operator: the spectral resolvent, applied to the true
  per-sector cascade profiles, produces a macroscopic sum converging to \(-\pi\)
  through collective mode weighting.
\item
  \textbf{Extend the bit-sharing model.} Section 9.5 shows the W-bit extended
  Markov chain reproduces Q\_emp ≈ 0.55, but R\^{}\{cas\} is unstable because
  σ\^{}\{cas\}₀₀ passes through zero. A better treatment of the entry carry at
  the window boundary, or an analytical argument showing q₀ → 1/2 as W → ∞,
  would connect the bit-sharing mechanism to the Diaconis--Fulman eigenvalue
  {[}P1, §8{]}.
\item
  \textbf{Proposition 2 (J = 1 exclusion): formal proof.} Verified
  computationally for K = 5, \ldots, 13. A formal case analysis on digit
  configurations near the boundary would promote this to a theorem.
\item
  \textbf{K ≥ 22 data.} Each additional K approximately quadruples the
  computation time. Data through K = 21 are available; K = 22 would require
  approximately 16 hours. Additional data points would sharpen the Richardson
  extrapolation and allow independent convergence analysis of the even-K and
  odd-K subsequences.
\item
  \textbf{Base dependence.} In base b ≥ 3, the sector ratio appears rational
  (e.g., R₅ = 5/4 for base 5). The angular uniqueness of base 2 is proved in
  {[}G{]}: the D-parity boundary is a straight line in angular coordinates iff b
  = 2, which explains why π enters only for binary multiplication. A systematic
  exact enumeration of R(K; b) for bases 3, 5, 7 would test whether R\_b is
  rational for all b ≥ 3.
\item
  \textbf{Boundary layer scaling.} Section 9.4 reports the TV(j) profile and
  c\_\{M−1\} distributions at a single K value (K = 12). Computing these at K =
  14, 16, 18 would reveal whether the key percentages (30\% vs 5.5\% for val =
  +1) converge to limiting values and at what rate.
\item
  \textbf{Separate weight sequences.} The current pipeline computes only the
  ratio R = σ₁₀/σ₀₀. Computing ⟨val⟩₁₀ and ⟨val⟩₀₀ separately would allow
  independent Richardson extrapolation and closed-form identification of each
  weight, which would provide a more transparent path to the identity
  ⟨val⟩₁₀/⟨val⟩₀₀ = −9.8197\ldots.
\end{enumerate}

\begin{center}\rule{0.5\linewidth}{0.5pt}\end{center}

\newpage

\subsection{References}\label{references}

\begin{enumerate}
\def\labelenumi{\arabic{enumi}.}
\tightlist
\item
  P. Diaconis and J. Fulman, \emph{Carries, shuffling, and symmetric functions},
  Adv. Appl. Math. \textbf{43} (2009), 176--196.
\item
  P. Diaconis and J. Fulman, \emph{Carries, shuffling, and an amazing matrix},
  Amer. Math. Monthly \textbf{116} (2009), 788--803.
\item
  J. Fulman, \emph{The carries process revisited}, preprint (2023).
\item
  J. M. Holte, \emph{Carries, combinatorics, and an amazing matrix}, Amer. Math.
  Monthly \textbf{104} (1997), 138--149.
\item
  D. E. Knuth, \emph{The Art of Computer Programming, Volume 2: Seminumerical
  Algorithms}, 3rd ed., Addison-Wesley, 1997.
\item
  L. F. Richardson and J. A. Gaunt, \emph{The deferred approach to the limit},
  Phil. Trans. R. Soc. A \textbf{226} (1927), 299--361.
\item
  {[}P1{]} Companion paper: \emph{π from Pure Arithmetic: A Spectral Phase
  Transition in the Binary Carry Bridge}. doi:10.5281/zenodo.18895611
\item
  {[}G{]} Companion paper: \emph{The angular uniqueness of base 2 in positional
  multiplication}. doi:10.5281/zenodo.18895601
\item
  T. Apostol, \emph{Introduction to Analytic Number Theory}, Springer, 1976.
\item
  {[}E{]} Companion paper: \emph{The Trace Anomaly of Binary Multiplication}.
  (Identifies conditionally the mechanism for \(R \to -\pi\) via resolvent
  universality, assuming the LMH.). doi:10.5281/zenodo.18895604
\end{enumerate}

\end{document}
